% GENERATED FILE. Edit canonical lessons, then run npm run generate:books.
% canonical-source-hash: fnv1a64:43818738d89e0361
% canonical-lessons: MR-C500-sangopan, MR-C500-sambhalne, MR-C500-jopasne, MR-C500-lad-karne, MR-C500-dattak-ghene

\chapter{Raising Children}
\label{ch:mr-raising-children}
\hlchaptermodality{500}

\begin{chapteropening}
\textbf{By the end of this chapter:} \emph{I can talk about bringing up, looking after, pampering and adopting a child.}

Complete the last lesson of chapter 500: I can talk about bringing up, looking after, pampering and adopting a child.
\end{chapteropening}

\section[\texorpdfstring{\mr{संगोपन} (saṅgopan)}{संगोपन (saṅgopan)}]{\mr{संगोपन} (saṅgopan) — upbringing}

\label{lesson:MR-C500-sangopan}



\begin{quote}
Before the new one: say the Marathi for to crawl, then the Marathi for to give to drink.
\end{quote}

\subsection*{You'll want to know: \mr{संगोपन}}
\textbf{\mr{संगोपन}} — \emph{saṅgopan} — \textquotedblleft{}upbringing\textquotedblright{}.

\begin{quote}
\mr{आजीने} \mr{माझं} \mr{संगोपन} \mr{केलं}.
\end{quote}

\begin{quote}
\emph{ājīne mājhaṁ saṅgopan kelaṁ.}
\end{quote}

\begin{quote}
My grandmother brought me up.
\end{quote}

Neuter.

\mr{संगोपन} \mr{करणे}, to bring up.

\begin{cousinweb}
\textbf{\mr{संगोपन}} (\emph{saṅgopan}) is a tatsama, taken from Sanskrit \emph{saṃgopana}, \textquotedblleft{}hiding or guarding well\textquotedblright{} (Monier-Williams); in Marathi it is the guarding of a child as it grows.
\end{cousinweb}

\begin{grammarlens}[title={What you've built}]
One more word for describing what you lived and felt.
\end{grammarlens}

\subsection*{Your turn}
\begin{itemize}
\raggedright
  \item \emph{Say it:} \emph{saṅgopan}
  \item \emph{Say it:} \emph{saṅgopan}, once more
  \item \emph{Say it:} \emph{ājīne mājhaṁ saṅgopan kelaṁ.} — My grandmother brought me up.
  \item \emph{Recall:} say the Marathi for to crawl, then the Marathi for to give to drink, then say \emph{saṅgopan} again
\end{itemize}

\begin{tcolorbox}[breakable,colback=teal!4,colframe=teal!35!black,title={Before you move on}]
What does \mr{संगोपन} mean? (\textquotedblleft{}Upbringing\textquotedblright{}.) Say it once more.
\end{tcolorbox}
\section[\texorpdfstring{\mr{सांभाळणे} (sāmbhāḷṇe)}{सांभाळणे (sāmbhāḷṇe)}]{\mr{सांभाळणे} (sāmbhāḷṇe) — to take care of; to look after}

\label{lesson:MR-C500-sambhalne}



\begin{quote}
Before the new one: say the Marathi for to give to drink, then the Marathi for upbringing.
\end{quote}

\subsection*{You'll want to know: \mr{सांभाळणे}}
\textbf{\mr{सांभाळणे}} — \emph{sāmbhāḷṇe} — \textquotedblleft{}to take care of; to look after\textquotedblright{}.

\begin{quote}
\mr{आई} \mr{कामावर} \mr{गेली} \mr{की} \mr{आजी} \mr{बाळाला} \mr{सांभाळते}.
\end{quote}

\begin{quote}
\emph{āī kāmāvar gelī kī ājī bāḷālā sāmbhāḷte.}
\end{quote}

\begin{quote}
When Mother goes to work, Grandmother looks after the baby.
\end{quote}

Past, with -\mr{ने} (he …): \mr{त्याने} \mr{बाळ} \mr{सांभाळलं}, \emph{tyāne bāḷ sāmbhāḷlaṁ}.

Also to handle with care: \mr{सांभाळून}, carefully.

\begin{cousinweb}
\textbf{\mr{सांभाळणे}} (\emph{sāmbhāḷṇe}) is inherited: CDIAL lists Marathi \emph{sā̃bhāḷṇẽ}, \textquotedblleft{}to take care of\textquotedblright{}, under Sanskrit \emph{sambharati} (CDIAL 12961).
\end{cousinweb}

\begin{grammarlens}[title={What you've built}]
One more word for describing what you lived and felt.
\end{grammarlens}

\subsection*{Your turn}
\begin{itemize}
\raggedright
  \item \emph{Say it:} \emph{sāmbhāḷṇe}
  \item \emph{Say it:} \emph{sāmbhāḷṇe}, once more
  \item \emph{Say it:} \emph{āī kāmāvar gelī kī ājī bāḷālā sāmbhāḷte.} — When Mother goes to work, Grandmother looks after the baby.
  \item \emph{Recall:} say the Marathi for to give to drink, then the Marathi for upbringing, then say \emph{sāmbhāḷṇe} again
\end{itemize}

\begin{tcolorbox}[breakable,colback=teal!4,colframe=teal!35!black,title={Before you move on}]
What does \mr{सांभाळणे} mean? (\textquotedblleft{}To take care of; to look after\textquotedblright{}.) Say it once more.
\end{tcolorbox}
\section[\texorpdfstring{\mr{जोपासणे} (jopāsṇe)}{जोपासणे (jopāsṇe)}]{\mr{जोपासणे} (jopāsṇe) — to nurture}

\label{lesson:MR-C500-jopasne}



\begin{quote}
Before the new one: say the Marathi for upbringing, then the Marathi for to take care of.
\end{quote}

\subsection*{You'll want to know: \mr{जोपासणे}}
\textbf{\mr{जोपासणे}} — \emph{jopāsṇe} — \textquotedblleft{}to nurture\textquotedblright{}.

\begin{quote}
\mr{त्याने} \mr{अनेक} \mr{वर्षं} \mr{ही} \mr{बाग} \mr{जोपासली}.
\end{quote}

\begin{quote}
\emph{tyāne anek varṣaṁ hī bāg jopāslī.}
\end{quote}

\begin{quote}
He nurtured this garden for many years.
\end{quote}

Past, with -\mr{ने} (he …): \mr{त्याने} \mr{बाग} \mr{जोपासली}, \emph{tyāne bāg jopāslī}.

Of a plant, a talent or a friendship.

\begin{cousinweb}
\textbf{\mr{जोपासणे}} (\emph{jopāsṇe}) is a Marathi verb whose origin is not settled.
\end{cousinweb}

\begin{grammarlens}[title={What you've built}]
One more word for describing what you lived and felt.
\end{grammarlens}

\subsection*{Your turn}
\begin{itemize}
\raggedright
  \item \emph{Say it:} \emph{jopāsṇe}
  \item \emph{Say it:} \emph{jopāsṇe}, once more
  \item \emph{Say it:} \emph{tyāne anek varṣaṁ hī bāg jopāslī.} — He nurtured this garden for many years.
  \item \emph{Recall:} say the Marathi for upbringing, then the Marathi for to take care of, then say \emph{jopāsṇe} again
\end{itemize}

\begin{tcolorbox}[breakable,colback=teal!4,colframe=teal!35!black,title={Before you move on}]
What does \mr{जोपासणे} mean? (\textquotedblleft{}To nurture\textquotedblright{}.) Say it once more.
\end{tcolorbox}
\section[\texorpdfstring{\mr{लाड} \mr{करणे} (lāḍ karṇe)}{लाड करणे (lāḍ karṇe)}]{\mr{लाड} \mr{करणे} (lāḍ karṇe) — to pamper}

\label{lesson:MR-C500-lad-karne}



\begin{quote}
Before the new one: say the Marathi for to take care of, then the Marathi for to nurture.
\end{quote}

\subsection*{You'll want to know: \mr{लाड} \mr{करणे}}
\textbf{\mr{लाड} \mr{करणे}} — \emph{lāḍ karṇe} — \textquotedblleft{}to pamper\textquotedblright{}.

\begin{quote}
\mr{आजोबा} \mr{नातवाचे} \mr{खूप} \mr{लाड} \mr{करतात}.
\end{quote}

\begin{quote}
\emph{ājobā nātvāce khūp lāḍ kartāt.}
\end{quote}

\begin{quote}
Grandfather pampers his grandson a lot.
\end{quote}

Past, with -\mr{ने} (he …): \mr{त्याने} \mr{लाड} \mr{केले}, \emph{tyāne lāḍ kele}.

\mr{लाड} is plural here; the one pampered takes the genitive: \mr{मुलाचे} \mr{लाड} \mr{करणे}.

\begin{cousinweb}
\textbf{\mr{लाड} \mr{करणे}} (\emph{lāḍ karṇe}) is a Marathi phrase built on \mr{लाड}, \textquotedblleft{}fondling\textquotedblright{}, which CDIAL lists under the reconstructed \emph{*lāḍya} (CDIAL 11013), with \mr{करणे}.
\end{cousinweb}

\begin{grammarlens}[title={What you've built}]
One more word for describing what you lived and felt.
\end{grammarlens}

\subsection*{Your turn}
\begin{itemize}
\raggedright
  \item \emph{Say it:} \emph{lāḍ karṇe}
  \item \emph{Say it:} \emph{lāḍ karṇe}, once more
  \item \emph{Say it:} \emph{ājobā nātvāce khūp lāḍ kartāt.} — Grandfather pampers his grandson a lot.
  \item \emph{Recall:} say the Marathi for to take care of, then the Marathi for to nurture, then say \emph{lāḍ karṇe} again
\end{itemize}

\begin{tcolorbox}[breakable,colback=teal!4,colframe=teal!35!black,title={Before you move on}]
What does \mr{लाड} \mr{करणे} mean? (\textquotedblleft{}To pamper\textquotedblright{}.) Say it once more.
\end{tcolorbox}
\section[\texorpdfstring{\mr{दत्तक} \mr{घेणे} (dattak gheṇe)}{दत्तक घेणे (dattak gheṇe)}]{\mr{दत्तक} \mr{घेणे} (dattak gheṇe) — to adopt}

\label{lesson:MR-C500-dattak-ghene}



\begin{quote}
Before the new one: say the Marathi for to nurture, then the Marathi for to pamper.
\end{quote}

\subsection*{You'll want to know: \mr{दत्तक} \mr{घेणे}}
\textbf{\mr{दत्तक} \mr{घेणे}} — \emph{dattak gheṇe} — \textquotedblleft{}to adopt\textquotedblright{}.

\begin{quote}
\mr{त्यांनी} \mr{एक} \mr{मुलगी} \mr{दत्तक} \mr{घेतली}.
\end{quote}

\begin{quote}
\emph{tyānnī ek mulgī dattak ghetlī.}
\end{quote}

\begin{quote}
They adopted a girl.
\end{quote}

Past, with -\mr{ने} (he …): \mr{त्याने} \mr{बाळ} \mr{दत्तक} \mr{घेतलं}, \emph{tyāne bāḷ dattak ghetlaṁ}.

\mr{दत्तक}, adopted: \mr{दत्तक} \mr{मुलगा}, an adopted son.

\begin{cousinweb}
\textbf{\mr{दत्तक} \mr{घेणे}} (\emph{dattak gheṇe}) is a Marathi phrase built on \mr{दत्तक}, a tatsama of Sanskrit \emph{dattaka}, a son given in adoption (Monier-Williams), with \mr{घेणे}, \textquotedblleft{}to take\textquotedblright{}.
\end{cousinweb}

\begin{grammarlens}[title={What you've built}]
One more word for describing what you lived and felt.
\end{grammarlens}

\subsection*{Your turn}
\begin{itemize}
\raggedright
  \item \emph{Say it:} \emph{dattak gheṇe}
  \item \emph{Say it:} \emph{dattak gheṇe}, once more
  \item \emph{Say it:} \emph{tyānnī ek mulgī dattak ghetlī.} — They adopted a girl.
  \item \emph{Recall:} say the Marathi for to nurture, then the Marathi for to pamper, then say \emph{dattak gheṇe} again
\end{itemize}

\begin{tcolorbox}[breakable,colback=teal!4,colframe=teal!35!black,title={Before you move on}]
What does \mr{दत्तक} \mr{घेणे} mean? (\textquotedblleft{}To adopt\textquotedblright{}.) Say it once more.
\end{tcolorbox}
